\section{Conclusion}
\label{sec:conclusion}

We presented MAW, an end-to-end framework that develops multi-agent orchestration by synthesizing multi-agent environments and training an orchestrator inside them.
Recursive environment improvement composes the reference graphs of single-agent tasks, reconstructs each composed graph into one coherent request, and admits it only after execution, replay, and semantic review.
The same construction verifies which subtasks are independent, which is what a schedule-aware process reward needs in order to score the whole team.
MAW thereby derives the environment in which the agents act and the ground truth that scores them from one construction.
Experiments on six agentic benchmarks show that this design improves task success and shortens the execution critical path.

The framework has four limitations.
Reward matching is anchored on a reference execution, which gives a correct alternative solution less process credit than a reference-matching one.
The schedule term also covers only the tasks whose independence survives a reordered execution.
Measured in logical model rounds, the critical path omits tool latency, so we establish a shorter critical path and not a proportional reduction in wall-clock time.
Requiring at least one delegation, the current interface never teaches the orchestrator when to solve a task alone, and evidence from one model scale with frozen workers leaves jointly trained teams for future work.
